\section{硬件实现思路}

\subsection{总体思路}

硬件不采用现成开发板，而是围绕 ESP32-S3-WROOM-1-N16R8 模组自行配置控制板。该模组包含 16\,MB Flash 和 8\,MB PSRAM，可满足固件、模型参数和低分辨率图像缓冲的需求。控制板集成 3.3\,V 电源、复位与启动电路、USB-C 接口、PCA9685 PWM 控制电路，并为 OV2640、MPU6050、OLED 和五路 MG90S 预留连接器。

ESP32-S3 模组由 3.3\,V 电源轨供电。根据模组数据手册，外部电源供电能力不应低于 0.5\,A；考虑摄像头和其他外设后，控制板的 3.3\,V 稳压电路按不低于 1\,A 的连续输出能力设计，并在模组和摄像头附近布置去耦电容。USB-C 的 D-、D+ 分别连接 GPIO19、GPIO20，用于下载、调试和上位机通信，接口处按硬件设计规范配置 CC 电阻、串联电阻和静电防护。

五路 MG90S 由独立的 5\,V 电源轨供电，PCA9685 只负责产生控制 PWM。PCA9685 的 VDD 接 3.3\,V，PWM 频率设为 50\,Hz，在 12 位计数下单步约为 4.88\,$\mu$s，实际脉宽与关节角的关系仍需逐轴标定。五路舵机同时启动时可能产生安培量级的瞬态电流，5\,V 电源按实测峰值留出余量，并在舵机电源入口设置大容量电解电容和就近陶瓷电容。逻辑域与舵机电源域共地，但舵机回流路径与数字电源分开布置。PCA9685 的 \texttt{OE} 引脚采用上拉默认关闭，待固件完成初始化后再使能输出。若 3.3\,V PWM 高电平在实测中不能可靠驱动 MG90S，则在信号端增加 74HCT 系列缓冲器转换到 5\,V。

\subsection{接口与资源约束}

OV2640 的并行接口需要 8 位数据线、像素时钟、场同步、行参考和外部时钟，另需 SCCB 配置线；OLED 使用 SPI，PCA9685 与 MPU6050 可共享 I²C。原理图冻结前需完成 GPIO 分配表，保留 GPIO19/20 给原生 USB，并避开启动配置脚以及模组内部 Flash/PSRAM 占用的引脚。摄像头连接器与 ESP32-S3 之间的并行信号应尽量短，并为时钟线预留串联阻尼位置。首版固件统一采用 ESP-IDF，并锁定与所选 ESP-WHO 版本兼容的 ESP-IDF 版本和提交号，以便复现实验环境。

\subsection{关节配置}

系统共有四个运动关节和一个夹爪执行轴，五路均使用 MG90S 位置舵机，配置见表~\ref{tab:servo-selection}。

\begin{table}[h]
\centering
\caption{舵机选型}
\label{tab:servo-selection}
\small
\begin{tabularx}{\linewidth}{l c Y}
\toprule
\textbf{关节} & \textbf{型号} & \textbf{主要设计约束}\\
\midrule
底座旋转 & MG90S & 配置独立轴承或转盘承重，舵机只提供旋转转矩\\
肩关节 & MG90S & 承担主要重力矩，需要重点核验力矩余量和温升\\
肘关节 & MG90S & 根据远端质量核验静力矩，并限制最大臂展负载\\
腕关节 & MG90S & 减轻夹爪与连接件质量，降低远端惯量\\
夹爪 & MG90S & 通过连杆驱动开合，并设置机械行程限制\\
\bottomrule
\end{tabularx}
\end{table}

\Needspace{7\baselineskip}
不同批次 MG90S 的实际扭矩、死区和回差可能存在差异，因此不能只依据商品页标称值判断是否满足负载要求。结构尺寸冻结前，肩、肘关节按
\[
  \tau_j=\sum_i m_i g r_{ij}
\]
估算各关节在最大伸展姿态下的静力矩，其中 $r_{ij}$ 为第 $i$ 个远端部件质心到关节 $j$ 的力臂，并至少保留约 2 倍余量；装配后再测试温升、保持能力和启动电流。

\subsection{传感器和外设}

\begin{itemize}
  \item \textbf{OV2640 摄像头}：通过控制板上的专用连接器接入 DVP/SCCB 信号。初期使用 QQVGA（$160\times120$）图像和单帧缓冲，以降低存储和带宽压力；实际帧率目标为 12--15\,FPS。摄像头安装在与固定底座刚性连接的支架上，不随机械臂关节运动，以保证桌面标定参数有效。
  \item \textbf{MPU6050 六轴传感器}：通过 I²C 连接，用于获取加速度和角速度。姿态可由 DMP 或主控上的互补滤波算法估计。无线手环方案需另配无线从节点，因为 MPU6050 本身不具备蓝牙功能；首版优先使用有线方案，减少无线任务对视觉联调的影响。
  \item \textbf{OLED 显示屏}：通过控制板上的 SPI 接口连接，显示当前模式、连接状态、目标是否丢失以及故障提示。
  \item \textbf{供电}：控制板逻辑部分使用 3.3\,V 电源轨，MG90S 使用独立 5\,V 电源轨。两路电源共地，但舵机大电流不经过 ESP32-S3 的稳压链路。
\end{itemize}

\subsection{机械结构}

机械结构采用串联形式，包括底座水平旋转、肩关节俯仰、肘关节俯仰、腕关节俯仰和夹爪开合。若将夹爪作为独立执行轴计数，系统共有五路执行轴；就末端空间位姿而言，机械臂本体提供四个运动自由度。大臂长度初定为 8--10\,cm，前臂长度为 7--9\,cm，预期工作区是底座前方半径约 15\,cm 的扇形区域。

结构件使用3D打印，自配置控制板和 OLED 固定在底座平台上。底座旋转轴增加独立轴承或转盘结构，避免由 MG90S 输出轴直接承受整臂的轴向和弯曲载荷。舵机电源和配重布置在底部，底座尺寸依据最大伸展姿态下的倾覆力矩确定。样机阶段优先通过调整臂长、壁厚和配重解决刚度与稳定性问题，不在首版中引入新的电机驱动形式。
